% =====================================================================
%  Estándares y buenas prácticas de LaTeX — guía de referencia
%  Suite de documentos LaTeX · Nivel 0 (Meta/referencia) · clase article.
%  Compilar: latexmk -pdf estandares.tex
% =====================================================================
\documentclass[11pt,a4paper]{article}

% --- Base estándar (este documento es ejemplo de sí mismo) ---
\usepackage[utf8]{inputenc}
\usepackage[T1]{fontenc}
\usepackage{lmodern}
\usepackage{microtype}
\usepackage[spanish,es-noshorthands]{babel}
\usepackage{csquotes}
\usepackage{amsmath}
\usepackage{amssymb}
\usepackage{booktabs}
\usepackage{tabularx}
\usepackage[margin=2.3cm]{geometry}
\usepackage{enumitem}
\usepackage{xcolor}
\usepackage{listings}
\usepackage{hyperref}
\hypersetup{unicode, pdftitle={Estándares y buenas prácticas de LaTeX},
            pdfauthor={Rodrigo Martín Vidal Galán},
            colorlinks=true, linkcolor=blue!50!black, urlcolor=blue!50!black}
\usepackage[spanish,capitalise,nameinlink]{cleveref}

% --- Estilo de código ---
\definecolor{codeback}{RGB}{245,247,250}
\definecolor{codekw}{RGB}{20,60,120}
\definecolor{codecom}{RGB}{110,120,135}
\definecolor{okgreen}{RGB}{0,120,60}
\definecolor{badred}{RGB}{170,30,40}
\lstset{
  basicstyle=\ttfamily\footnotesize,
  commentstyle=\color{codecom}\itshape,
  backgroundcolor=\color{codeback},
  frame=single, rulecolor=\color{codekw!30},
  breaklines=true, showstringspaces=false, columns=fullflexible,
  xleftmargin=8pt, xrightmargin=4pt, framesep=4pt,
  extendedchars=true,
  literate={á}{{\'a}}1 {é}{{\'e}}1 {í}{{\'i}}1 {ó}{{\'o}}1 {ú}{{\'u}}1
           {ñ}{{\~n}}1 {¿}{{?`}}1 {¡}{{!`}}1 {°}{{\textdegree}}1,
}
\lstdefinestyle{tex}{language=[LaTeX]TeX, keywordstyle=\color{codekw}\bfseries}
\lstdefinestyle{sh}{language=bash, keywordstyle=\ttfamily}
\lstset{style=tex}
\newcommand{\C}[1]{\texttt{\textbackslash #1}}
\newcommand{\pkg}[1]{\textsf{#1}}
\newcommand{\si}{\textcolor{okgreen}{\boldmath$\checkmark$}}
\newcommand{\no}{\textcolor{badred}{\boldmath$\times$}}

\newenvironment{nota}{\par\smallskip\noindent\begin{center}%
  \begin{minipage}{0.95\linewidth}\small\rule{\linewidth}{0.4pt}\par\nobreak}%
  {\par\rule{\linewidth}{0.4pt}\end{minipage}\end{center}\smallskip}

\title{\textbf{Estándares y buenas prácticas de \LaTeX}\\[2pt]
       \large Guía de referencia}
\author{Rodrigo Martín Vidal Galán}
\date{Junio de 2026}

\begin{document}
\maketitle

\noindent Referencia para \enquote{hacer las cosas bien}: qué cargar, cómo
estructurar y qué evitar. Cada recomendación viene con su \textbf{porqué}. Pensá
este documento en tres niveles: la \textbf{base mínima} (lo que casi todo
documento debería tener), lo \textbf{recomendado} (el superconjunto razonable) y
los \textbf{anti-patrones} (lo que conviene no usar). Las bases listas para
copiar están en \cref{sec:bases}. \emph{(Para instalar y compilar, ver
\emph{Flujo e instalación}; para el universo de temas, \emph{TEMARIO}.)}

\newpage
\tableofcontents
\newpage

% =====================================================================
\section{Principios generales}
% =====================================================================
Toda la guía se apoya en cinco principios; el resto son consecuencias:
\begin{enumerate}[noitemsep,leftmargin=*]
  \item \textbf{Separar contenido de presentación.} Marcá \emph{qué es} cada cosa
        (semántico), no \emph{cómo se ve} (visual): definí \C{conjunto\{\dots\}} en
        vez de escribir las llaves a mano. Así, un cambio de estilo se hace en un
        solo lugar.
  \item \textbf{DRY} (\emph{Don't Repeat Yourself}). Lo que se repite va en una
        macro o entorno. Menos errores, cambios centralizados.
  \item \textbf{Compilar limpio.} La meta es \textbf{cero errores y warnings}; un
        \texttt{Overfull \textbackslash hbox} o un \texttt{Font shape undefined}
        son señales de algo a corregir, no ruido.
  \item \textbf{Reproducibilidad.} Que cualquiera obtenga el mismo PDF: versiones,
        \texttt{latexmk}, control de versiones.
  \item \textbf{Legibilidad del fuente.} El \texttt{.tex} se lee y se mantiene:
        comentarios, indentación y \textbf{una idea por línea}.
\end{enumerate}

% =====================================================================
\section{Motor de compilación}
% =====================================================================
El \textbf{motor} es lo primero que define todo lo demás (sobre todo el bloque de
fuentes). Las tres opciones:

\begin{center}\small
\renewcommand{\arraystretch}{1.3}
\begin{tabularx}{\linewidth}{@{}l XXX@{}}
\toprule
 & \textbf{pdfLaTeX} & \textbf{XeLaTeX} & \textbf{LuaLaTeX} \\
\midrule
Entrada & UTF-8 (con \pkg{inputenc}) & UTF-8 nativo & UTF-8 nativo \\
Fuentes & Solo las de TeX & Del sistema (\pkg{fontspec}) & Del sistema (\pkg{fontspec}) \\
\pkg{microtype} & \textbf{Completo} & Parcial & Casi completo \\
Velocidad & \textbf{La más rápida} & Media & La más lenta \\
Scripting & No & No & \textbf{Sí (Lua)} \\
Madurez & \textbf{Máxima} & Alta & Alta (futuro) \\
\bottomrule
\end{tabularx}
\end{center}

\noindent\textbf{Por qué importa y qué elegir.} \texttt{pdfLaTeX} sigue siendo el
\emph{default robusto} para matemática: máxima compatibilidad de paquetes y
\pkg{microtype} completo. \texttt{LuaLaTeX} es el estándar \textbf{emergente};
elegilo si necesitás fuentes del sistema, Unicode avanzado o programar con Lua.
\texttt{XeLaTeX} es similar a Lua para fuentes, pero con menos futuro.

\begin{nota}
\textbf{Caveat.}\enspace Muchas plantillas de revistas/congresos están hechas para
pdfLaTeX. Si vas a enviar a una, revisá sus \emph{submission guidelines} antes de
cambiar de motor.
\end{nota}

% =====================================================================
\section{Clase de documento}
% =====================================================================
La \textbf{clase} (\C{documentclass}) fija la estructura disponible. Según el
propósito: \texttt{article} (artículos, TPs), \texttt{report} (con capítulos),
\texttt{book} (libros), \texttt{beamer} (presentaciones), \texttt{standalone}
(una figura suelta).

\noindent\textbf{Alternativa \enquote{premium}:} las clases \textbf{KOMA-Script}
(\texttt{scrartcl}/\texttt{scrreprt}/\texttt{scrbook}) reemplazan a
\texttt{article}/\texttt{report}/\texttt{book} con mejor diseño por defecto y
comandos consistentes (\C{setkomafont}, \C{KOMAoptions}). \texttt{memoir} es otra
clase potente \enquote{todo en uno}. Para documentos largos/serios, muchos las
prefieren.

% =====================================================================
\section{Codificación y fuentes}
% =====================================================================
Este bloque resuelve tres problemas concretos que, mal hechos, arruinan el PDF:
acentos rotos, fuentes borrosas y guionado incorrecto.

\paragraph{pdfLaTeX (lo más usado).}
\begin{lstlisting}
\usepackage[utf8]{inputenc}   % entrada UTF-8 (redundante desde 2018, pero explicito)
\usepackage[T1]{fontenc}      % codificacion de SALIDA: acentos, guionado, copy-paste
\usepackage{lmodern}          % Latin Modern: fuentes vectoriales (evita bitmaps)
\usepackage{microtype}        % protrusion + expansion: el "pulido" tipografico
\end{lstlisting}
\begin{itemize}[noitemsep,leftmargin=*]
  \item \textbf{\pkg{fontenc} \texttt{[T1]} es casi obligatorio en español:} sin
        él, las palabras con acento \textbf{no se guionan} (quedan feos huecos) y
        el copy-paste del PDF falla (copiás \enquote{funci} en vez de
        \enquote{función}). Define cómo se codifican los caracteres en la salida.
  \item \textbf{\pkg{lmodern}} reemplaza las fuentes por defecto (que pueden salir
        como \textbf{bitmaps} borrosos al hacer zoom) por versiones vectoriales.
        Además, \pkg{microtype} necesita fuentes Type\,1 como estas para funcionar.
  \item \textbf{\pkg{microtype}} hace ajustes finísimos (saca un poco de tinta
        hacia el margen, estira/encoge letras imperceptiblemente) que mejoran el
        gris de la página y reducen los \emph{overfull}. Casi gratis; cargalo
        siempre.
\end{itemize}

\paragraph{XeLaTeX / LuaLaTeX (alternativa moderna).} \textbf{No} se usan
\pkg{inputenc}/\pkg{fontenc}; en su lugar:
\begin{lstlisting}
\usepackage{fontspec}
\setmainfont{Latin Modern Roman}   % o cualquier fuente del sistema
\usepackage{unicode-math}
\setmathfont{Latin Modern Math}
\end{lstlisting}

% =====================================================================
\section{Idioma}
% =====================================================================
\begin{lstlisting}
\usepackage[spanish,es-noshorthands]{babel}
\end{lstlisting}
\pkg{babel} con \texttt{spanish} activa el guionado en español, traduce
\enquote{Figura/Cuadro/Índice} y pone \C{today} en español.

\noindent\textbf{Por qué \texttt{es-noshorthands}.} El español de \pkg{babel}
define \enquote{shorthands}: hace que \verb|"|, \verb|~| y \verb|.| tomen
significados especiales. Eso es cómodo para tipear, pero \textbf{rompe}
\pkg{tikz}, \pkg{siunitx} y las URLs. Desactivarlos con \texttt{es-noshorthands}
evita horas de errores misteriosos.

\noindent\pkg{csquotes} + \C{enquote\{\dots\}} da \textbf{comillas tipográficas}
correctas según el idioma, en lugar de escribir \verb|``...''| a mano.

% =====================================================================
\section{Matemática (el núcleo)}
% =====================================================================
\paragraph{Paquetes.}
\begin{lstlisting}
\usepackage{mathtools}  % superconjunto de amsmath: \coloneqq, \DeclarePairedDelimiter, fixes
\usepackage{amssymb}    % \mathbb, \mathcal, \leq, flechas...
\usepackage{amsthm}     % entornos theorem/definition/proof
\usepackage{siunitx}    % numeros y unidades: \num, \qty
\usepackage{bm}         % negrita matematica correcta (\bm{x}) para vectores
\end{lstlisting}
\textbf{Clave:} cargá \pkg{mathtools}, \textbf{no} solo \pkg{amsmath}: lo incluye
y además corrige bugs y agrega herramientas. \pkg{amsmath} es prácticamente
obligatorio (te da \texttt{align}, \C{text}, \C{DeclareMathOperator}, el espaciado
correcto).

\paragraph{Macros semánticas (alta prioridad).} Definir atajos hace el documento
más corto, consistente y fácil de cambiar:
\begin{lstlisting}
\newcommand{\R}{\mathbb{R}}   \newcommand{\N}{\mathbb{N}}
\DeclareMathOperator{\sen}{sen}                 % seno en espanol (redonda)
\DeclarePairedDelimiter{\abs}{\lvert}{\rvert}   % \abs{x} y \abs*{...} auto-escalado
\end{lstlisting}
Los \textbf{operadores} con \C{DeclareMathOperator} salen en redonda y con el
espaciado correcto; los \textbf{delimitadores} con \C{DeclarePairedDelimiter}
(de \pkg{mathtools}) escalan solos.

\paragraph{Do / Don't.}
\begin{center}\small
\renewcommand{\arraystretch}{1.3}
\begin{tabularx}{\linewidth}{@{}l l X@{}}
\toprule
\textbf{\no\ Evitar} & \textbf{\si\ Usar} & \textbf{Por qué} \\
\midrule
\verb|$$...$$| & \C{[}\,\dots\,\C{]} / \texttt{equation} & \verb|$$| es TeX plano: rompe el espaciado \\
\texttt{eqnarray} & \texttt{align} & \texttt{eqnarray} tiene el espaciado roto \\
\verb;|x|; a mano & \C{abs\{x\}} & escala y es semántico \\
$sen\,x$ (pelado) & \C{sen} (operador) & espaciado y redonda correctos \\
\C{bf}, \C{it} en math & \C{mathbf}, \C{bm} & los primeros son obsoletos \\
\texttt{...} & \C{dots} & elige \texttt{\textbackslash ldots}/\texttt{\textbackslash cdots} según contexto \\
\bottomrule
\end{tabularx}
\end{center}

\paragraph{\pkg{siunitx} (números y unidades).} Evita escribir a mano separadores
de miles y unidades, y mantiene todo consistente:
\begin{lstlisting}
\num{55000}            % 55 000
\qty{10}{\celsius}     % 10 grados C  (v3; en v2 era \SI)
\qty{250}{\litre\per\minute}
\end{lstlisting}
En \pkg{siunitx} v3 los comandos viejos \C{SI}/\C{si} están desaconsejados: usá
\C{qty}/\C{unit}.

% =====================================================================
\section{Tipografía y diseño de página}
% =====================================================================
\begin{center}\small
\renewcommand{\arraystretch}{1.3}
\begin{tabularx}{\linewidth}{@{}l X@{}}
\toprule
\textbf{Paquete} & \textbf{Para qué (y por qué)} \\
\midrule
\pkg{microtype} & Pulido tipográfico. Casi siempre. \\
\pkg{geometry} & Márgenes (\C{usepackage[margin=2.5cm]\{geometry\}}), en vez de tocar \C{textwidth} a mano. \\
\pkg{setspace} & Interlineado (\C{onehalfspacing}) para tesis/borradores. \\
\pkg{parskip} & Espacio entre párrafos en lugar de sangría (estilo \enquote{informe}). \\
\pkg{enumitem} & Personalizar listas (espaciado, etiquetas). Reemplaza hacks con \C{vspace}. \\
\pkg{titlesec} & Formato de títulos de sección (si KOMA no alcanza). \\
\pkg{fancyhdr} & Encabezados/pies. \\
\bottomrule
\end{tabularx}
\end{center}
\textbf{Canon:} márgenes con \pkg{geometry} (no redefinir \C{textwidth} a mano);
listas con \pkg{enumitem}; \textbf{no} forzar saltos con \C{\textbackslash} en
texto corrido (eso es para tablas y \texttt{align}).

% =====================================================================
\section{Gráficos y figuras}
% =====================================================================
\begin{lstlisting}
\usepackage{graphicx}     % \includegraphics
\usepackage{pgfplots}     % graficos de funciones/datos (carga tikz)
\pgfplotsset{compat=1.18} % FIJAR la version: reproducibilidad
\usepackage{caption}\usepackage{subcaption}
\usepackage{float}        % opcion [H]
\end{lstlisting}
\begin{itemize}[noitemsep,leftmargin=*]
  \item \textbf{Vectorial} siempre que se pueda (PDF/TikZ/pgfplots); \emph{raster}
        (PNG/JPG) solo para fotos. Lo vectorial no se pixela al hacer zoom.
  \item Toda figura va en un entorno \texttt{figure} con \C{caption}
        \textbf{y} \C{label} (el caption \emph{antes} del label).
  \item \textbf{Fijar \texttt{compat}} en \pkg{pgfplots} evita que el aspecto
        cambie entre versiones del paquete.
  \item Con muchos TikZ, la \textbf{externalización} cachea cada figura y acelera
        la compilación.
\end{itemize}

% =====================================================================
\section{Tablas}
% =====================================================================
La regla de oro es \pkg{booktabs}: tablas con solo \C{toprule}, \C{midrule} y
\C{bottomrule}, \textbf{sin líneas verticales} ni dobles. ¿Por qué? Las líneas
verticales y el exceso de reglas \textbf{ensucian} y dificultan la lectura; el
estilo \enquote{de libro} usa aire y reglas horizontales finas. Para columnas que
ajustan su ancho, \pkg{tabularx}; para alinear números por el decimal, la columna
\texttt{S} de \pkg{siunitx}; para tablas largas, \pkg{longtable}. \textbf{Nunca}
\verb|\hline\hline| ni \verb#|c|c|#.

% =====================================================================
\section{Referencias cruzadas y bibliografía}
% =====================================================================
\begin{lstlisting}
\usepackage{hyperref}                          % enlaces + metadatos (casi ULTIMO)
\usepackage[capitalise,nameinlink]{cleveref}   % \cref -> "Figura 3"; DESPUES de hyperref
\end{lstlisting}
\C{label}/\C{ref} siempre; \textbf{mejor \C{cref}}, que agrega el tipo
(\enquote{figura}, \enquote{ecuación}) y se adapta al idioma, así no escribís
\enquote{la figura} a mano. \pkg{hyperref} se carga casi al final (excepción:
\pkg{cleveref} va \emph{después}).

\noindent\textbf{Bibliografía:} el enfoque \textbf{moderno} es \pkg{biblatex} +
\textbf{biber} (flexible, internacionalizado, estilos \texttt{numeric}/
\texttt{authoryear}/\texttt{ieee}\dots). El clásico \pkg{natbib} + BibTeX todavía
se usa porque algunas revistas lo exigen.

% =====================================================================
\section{Color y código}
% =====================================================================
\textbf{Color} con \pkg{xcolor}: definí colores \textbf{nombrados} una vez
(\C{definecolor\{azulUAI\}\{RGB\}\{\dots\}}) y reutilizalos; no dispersás
\C{color[RGB]\{\dots\}}. \textbf{Accesibilidad:} contraste suficiente y no
codificar información \emph{solo} por color.

\noindent\textbf{Mostrar código:} \pkg{listings} (sin dependencias, lo que usa
esta suite) o \pkg{minted} (resaltado superior con Pygments, pero requiere
\texttt{-shell-escape}). En \pkg{beamer}, todo \texttt{frame} con código necesita
\texttt{[fragile]}.

% =====================================================================
\section{Organización del proyecto y del código}
% =====================================================================
\begin{itemize}[noitemsep,leftmargin=*]
  \item \textbf{Modularizar} documentos grandes: \C{input} (pega texto),
        \C{include} (capítulos + \C{includeonly}), o \pkg{subfiles} (compilar cada
        parte sola). El \textbf{preámbulo} en un archivo aparte y reutilizable.
  \item \textbf{Una oración por línea} en el fuente: los diffs de git quedan
        limpios y se ve exactamente qué cambió.
  \item \textbf{Macros DRY} para lo repetido:
\end{itemize}
\begin{lstlisting}
\newcommand{\cs}[1]{\texttt{\textbackslash#1}}  % \cs{align} en vez de \texttt{\textbackslash align}
\end{lstlisting}

% =====================================================================
\section{Beamer (presentaciones)}
% =====================================================================
\pkg{beamer} ya carga \pkg{hyperref}, \pkg{xcolor}, \pkg{amsmath}, \pkg{tikz}
(no hace falta recargarlos). Lo esencial:
\begin{itemize}[noitemsep,leftmargin=*]
  \item \texttt{[fragile]} en todo \texttt{frame} con código verbatim/listings.
  \item \textbf{Una idea por diapositiva}; bloques semánticos
        (\texttt{block}/\texttt{alertblock}/\texttt{exampleblock}).
  \item \textbf{Overlays} (\C{pause}, \texttt{<n->}) con moderación; \texttt{handout}
        para imprimir sin animaciones.
  \item Anchos de columnas + separación \textbf{$\le$} \C{textwidth} (si no, se
        van al margen).
  \item \C{hypersetup\{unicode\}} para que el índice/marcadores muestren bien los
        acentos y \enquote{¿}.
\end{itemize}

% =====================================================================
\section{Reproducibilidad, compilación y diagnóstico}
% =====================================================================
\textbf{\texttt{latexmk}} es el estándar para compilar (resuelve pasadas, biber,
índices). El objetivo es un \textbf{log limpio}: revisá \texttt{Overfull},
\texttt{Font shape}, \texttt{Reference/Citation undefined}.

\noindent\textbf{Paquetes de diagnóstico} (quitar en la versión final) que
\emph{automatizan} el control de calidad mientras escribís:
\begin{lstlisting}
\usepackage[l2tabu,orthodox]{nag}     % AVISA en el log si usas comandos obsoletos
\usepackage[notref,notcite]{showkeys} % muestra los \label en el PDF (borrador)
\overfullrule=2cm                     % marca con una barra los Overfull \hbox
\usepackage{todonotes}                % \todo{...} y \listoftodos
\end{lstlisting}
\pkg{nag} es el complemento natural de \pkg{l2tabu} (\cref{sec:anti}): convierte
la lista de anti-patrones en \textbf{avisos automáticos al compilar}.

% =====================================================================
\section{Anti-patrones (qué NO usar)}\label{sec:anti}
% =====================================================================
Resumen de \pkg{l2tabu} (\emph{An essential guide to \LaTeXe{} usage: Obsolete
commands and packages}). Los obsoletos suelen \textbf{romper} el sistema de
fuentes moderno o tener bugs:
\begin{center}\small
\renewcommand{\arraystretch}{1.25}
\begin{tabularx}{\linewidth}{@{}l X@{}}
\toprule
\textbf{\no\ Obsoleto / malo} & \textbf{\si\ Reemplazo} \\
\midrule
\C{bf} \C{it} \C{rm} \C{sc} \C{sl} \C{tt} & \C{textbf} \C{textit}\dots / \C{bfseries}\dots \\
\verb|$$...$$| & \C{[}\,\dots\,\C{]} \\
\texttt{eqnarray} & \texttt{align} \\
\C{over} & \C{frac} \\
\C{centerline} & \C{centering} / \texttt{center} \\
\texttt{a4}, \texttt{a4wide} & \pkg{geometry} \\
\pkg{epsfig}, \pkg{psfig} & \pkg{graphicx} \\
\C{def} para todo & \C{newcommand} (chequea redefiniciones) \\
Tablas con \verb#|# y \verb|\hline\hline| & \pkg{booktabs} \\
\bottomrule
\end{tabularx}
\end{center}

% =====================================================================
\section{Bases concretas: MÍNIMA y MÁXIMA}\label{sec:bases}
% =====================================================================
\textbf{\enquote{Máxima} no significa \enquote{cargar todo siempre}}, sino el
\emph{superconjunto recomendado}; en cada documento se usa el subconjunto que el
contenido pide (cargar paquetes que no se usan también es un anti-patrón).

\paragraph{\texttt{article} — base MÍNIMA.}
\begin{lstlisting}
\documentclass[11pt,a4paper]{article}
\usepackage[utf8]{inputenc}
\usepackage[T1]{fontenc}
\usepackage{lmodern}
\usepackage[spanish,es-noshorthands]{babel}
\usepackage{amsmath,amssymb}
\usepackage{microtype}
\usepackage{hyperref}
\end{lstlisting}

\paragraph{\texttt{article} — base MÁXIMA (matemática).}
\begin{lstlisting}
\documentclass[11pt,a4paper]{article}
\usepackage[utf8]{inputenc}\usepackage[T1]{fontenc}\usepackage{lmodern}\usepackage{microtype}
\usepackage[spanish,es-noshorthands]{babel}\usepackage{csquotes}
\usepackage{mathtools}\usepackage{amssymb}\usepackage{amsthm}\usepackage{bm}\usepackage{siunitx}
\usepackage{booktabs}\usepackage{tabularx}
\usepackage{graphicx}\usepackage{pgfplots}\pgfplotsset{compat=1.18}
\usepackage{caption}\usepackage{subcaption}\usepackage{float}
\usepackage[margin=2.5cm]{geometry}\usepackage{enumitem}
\usepackage[backend=biber,style=numeric]{biblatex}\addbibresource{refs.bib}
\usepackage{hyperref}\usepackage{cleveref}     % cleveref DESPUES de hyperref
\hypersetup{unicode, pdftitle={...}, pdfauthor={...}, hidelinks}
\newcommand{\R}{\mathbb{R}}\DeclareMathOperator{\sen}{sen}
\DeclarePairedDelimiter{\abs}{\lvert}{\rvert}
\theoremstyle{plain}\newtheorem{teo}{Teorema}[section]
\end{lstlisting}

\paragraph{\texttt{beamer} — base MÍNIMA.}
\begin{lstlisting}
\documentclass[10pt,aspectratio=169]{beamer}
\usetheme{Madrid}
\usepackage[utf8]{inputenc}\usepackage[T1]{fontenc}\usepackage{lmodern}
\usepackage[spanish,es-noshorthands]{babel}
\usepackage{amsmath,amssymb}
\hypersetup{unicode}   % beamer ya cargo hyperref
\end{lstlisting}

\paragraph{\texttt{beamer} — base MÁXIMA (presentación de matemática).}
\begin{lstlisting}
\documentclass[10pt,aspectratio=169]{beamer}
\usetheme{Madrid}\usecolortheme{whale}\setbeamertemplate{navigation symbols}{}
\usepackage[utf8]{inputenc}\usepackage[T1]{fontenc}\usepackage{lmodern}\usepackage{microtype}
\usepackage[spanish,es-noshorthands]{babel}\usepackage{csquotes}
\usepackage{mathtools}\usepackage{amssymb}\usepackage{bm}\usepackage{siunitx}
\usepackage{booktabs}\usepackage{pgfplots}\pgfplotsset{compat=1.18}
\usepackage{listings}
\hypersetup{unicode, pdftitle={...}, pdfauthor={...}}
\newcommand{\R}{\mathbb{R}}\DeclareMathOperator{\sen}{sen}
\end{lstlisting}

% =====================================================================
\section{Checklist antes de entregar}
% =====================================================================
\begin{itemize}[label=$\square$,leftmargin=*,itemsep=2pt]
  \item Compila \textbf{sin errores} y, idealmente, \textbf{sin warnings}.
  \item \pkg{fontenc}\texttt{[T1]} + \pkg{lmodern} (o \pkg{fontspec}); \pkg{microtype}.
  \item Idioma con \pkg{babel} y \texttt{es-noshorthands}.
  \item Matemática con \pkg{amsmath}/\pkg{mathtools}; nada de \verb|$$|/\texttt{eqnarray}.
  \item Figuras y tablas con \C{caption}+\C{label}; tablas con \pkg{booktabs}.
  \item Referencias con \C{cref}/\C{eqref}; sin \enquote{??} en el PDF.
  \item Metadatos del PDF (\texttt{pdftitle}, \texttt{pdfauthor}); \C{hypersetup\{unicode\}}.
  \item \texttt{.gitignore} para intermedios.
  \item (Beamer) \texttt{[fragile]} en frames con código.
\end{itemize}

% =====================================================================
\section{Validación contra fuentes (junio 2026)}
% =====================================================================
Esta guía se contrastó con fuentes web frescas: la base se \textbf{confirma}.
Confirmado: \pkg{inputenc} redundante desde 2018 (pero inofensivo); usar
\pkg{mathtools}; pdfLaTeX como default y LuaLaTeX como emergente;
\pkg{fontenc}+\pkg{lmodern}+\pkg{microtype}; \pkg{biblatex}+biber moderno.
Sutilezas incorporadas: \pkg{siunitx} v3 (\C{SI}$\to$\C{qty}); en Beamer
\texttt{metropolis} está sin mantenimiento y su fork \texttt{moloch} lo continúa;
los \textbf{paquetes de diagnóstico} (\pkg{nag}\dots); y opciones
\texttt{cleveref[capitalise,nameinlink]}, \texttt{csquotes[strict=true]}.

\paragraph{Recursos autoritativos.}
\begin{itemize}[noitemsep,leftmargin=*]
  \item \pkg{l2tabu} — comandos/paquetes obsoletos: \url{https://ctan.org/pkg/l2tabu}
  \item \emph{The Not So Short Introduction to LaTeX} (\texttt{lshort}).
  \item \url{https://tex.stackexchange.com} y \texttt{texdoc <paquete>}.
  \item Manuales: \texttt{texdoc amsmath}, \texttt{mathtools}, \texttt{siunitx},
        \texttt{biblatex}, \texttt{beamer}, \texttt{microtype}.
\end{itemize}

% =====================================================================
\section*{Ver también}
% =====================================================================
\addcontentsline{toc}{section}{Ver también}
\begin{itemize}[noitemsep,leftmargin=*]
  \item \emph{Flujo e instalación} — compilar, distribuciones, herramientas.
  \item \emph{Plantillas} — \texttt{article.tex} y \texttt{beamer.tex} con estas bases.
  \item \emph{TEMARIO} y \emph{README} de la suite.
\end{itemize}

\end{document}
